% problem_id: S014-lurie-ultracategories-proposition-4-2-3
% source_id: S014 (Jacob Lurie, Ultracategories)
% source_file: lurie-conceptual.pdf
% statement_locator: printed p. 47
% proof_locator: printed p. [50]
% representation: PDF; transcribed from rendered page images (never from the text layer)
% status: complete (statement + author proof verbatim + reason + locators)
%
%% problem_id: S014-lurie-ultracategories-prop-4.2.3
%% source_id: S014 (Jacob Lurie, Ultracategories)
%% source_file: lurie-conceptual.pdf  (host: /disks/sata1/yupeng/human-proof-corpus/source-cache/registry-sources/lurie/lurie-conceptual.pdf)
%% candidate: M2
%% target item: Proposition 4.2.3, section 4.2 ("Free Objects of Comp_M")
%% statement_locator: printed page 47 = PDF page 47 (page header "ULTRACATEGORIES   47")
%% proof_locator: printed page 50 = PDF page 50 (page header "50   ULTRACATEGORIES")
%% representation: PDF; transcribed by eye from rendered page images (pdftoppm -r 150 -png for whole pages,
%%                 plus -r 400 -png crops for every symbol-level check). The PDF text layer is corrupt
%%                 (it renders "\in" as "2", "\alpha" as a dead glyph, drops "\theta"/"\epsilon" entirely
%%                 and loses spaces around reference numbers), so it was used ONLY to locate line numbers;
%%                 no character below was taken from the text layer.
%% status: TRANSCRIPTION ONLY (difficulty judgement recorded separately in M2.json)
%%
%% Transcription notes:
%%  * Pages transcribed, in page order: 47, 48, 49, 50. Beside Proposition 4.2.3 itself, the surrounding
%%    items of section 4.2 are reproduced (Question 4.2.1, Remarks 4.2.2/4.2.4/4.2.6, Notation 4.2.5,
%%    Theorem 4.2.7, Propositions 4.2.8/4.2.9/4.2.10 with their proofs), because the proof of Proposition
%%    4.2.3 refers to Propositions 4.2.8 and 4.2.9 and those are the items that carry its content. Nothing
%%    outside pages 47--50 is transcribed.
%%  * Boxed cross-reference numbers of the source (hyperref boxes) are transcribed as plain references.
%%  * The top of printed page 47 holds the tail of the proof of a Proposition of section 4.1 (the passage
%%    computing O_{X,int_S fbar(s)dmu}); it does not belong to section 4.2 and is NOT transcribed here.
%%  * The "Proof of Theorem 4.2.7" begins at the foot of printed page 50 and continues onto printed page 51,
%%    which was not rendered; only its opening sentence is reproduced, and the cut is marked.
%%  * Places where the printed text deviates from what the mathematics requires are kept EXACTLY as printed
%%    and flagged by a %% [NOTE: ...] line. In particular the source intermittently prints "N_t" where the
%%    surrounding formulas require the object family "M_t"; this was verified against a 400 dpi crop in which
%%    both glyphs occur side by side in one formula (see the note at the condition-(2) diagram below).
%%    Nothing has been silently corrected.
%%  * Commutative diagrams are transcribed into "array" environments: every node, arrow and arrow label is
%%    reproduced verbatim, the geometric layout is approximated (marked by a %% [NOTE: ...] line each time).
%%  * No passage of the transcribed pages was illegible, so there are no [ILLEGIBLE: ...] markers.

\documentclass[11pt]{article}
\usepackage{amsmath}
\usepackage{amssymb}
\usepackage{amsthm}
\usepackage{mathrsfs}
\usepackage[margin=1in]{geometry}

\theoremstyle{plain}
\newtheorem{theorem}{Theorem}
\newtheorem{proposition}{Proposition}
\newtheorem{lemma}{Lemma}
\newtheorem{corollary}{Corollary}
\newtheorem{question}{Question}

\theoremstyle{definition}
\newtheorem{remark}{Remark}
\newtheorem{notation}{Notation}

\newcommand{\Comp}{\operatorname{Comp}}
\newcommand{\FunLUlt}{\operatorname{Fun}^{\mathrm{LUlt}}}
\newcommand{\Hom}{\operatorname{Hom}}

\begin{document}

%% --- page 47 ---

%% [NOTE: printed page 47 opens with the tail of a proof belonging to section 4.1; not transcribed, see the
%%  transcription notes above.]

\subsection*{4.2. Free Objects of $\Comp_{\mathcal{M}}$.}

In \S 1.3, we introduced the notion of an \emph{ultrastructure} on a category $\mathcal{M}$ (Definition 1.3.1). Our definition of ultrastructure was intended to axiomatize the essential features of the categorical ultraproducts studied in \S 1.2: whenever $\mathcal{M}$ has ultraproducts in some larger category $\mathcal{M}^{+}$, the formation of ultraproducts determines an ultrastructure on $\mathcal{M}$ (Proposition 1.3.7). It is natural to ask the following:

\begin{question}[4.2.1]
Let $\mathcal{M}$ be a category. Can every ultrastructure on $\mathcal{M}$ be obtained by embedding $\mathcal{M}$ into some larger category $\mathcal{M}^{+}$, such that $\mathcal{M}$ has ultraproducts in $\mathcal{M}^{+}$?
\end{question}

In this section, we give an affirmative answer to Question 4.2.1. To every ultracategory $\mathcal{M}$, we show that there is a \emph{canonical} way to recover the ultrastructure via categorical ultraproducts in a larger category $\mathcal{M}^{+}$: namely, we can take $\mathcal{M}^{+}$ to be the opposite of the category $\Comp_{\mathcal{M}}$ of Construction 4.1.1 (Theorem 4.2.7).

\begin{remark}[4.2.2]
For every ultracategory $\mathcal{M}$, there exists an embedding $\mathcal{M} \hookrightarrow \mathcal{M}^{+}$ satisfying the requirement of Question 4.2.1. However, the enlargement $\mathcal{M}^{+}$ is not uniquely determined. Moreover, the construction of this section is not the most efficient: to recover the ultrastructure on $\mathcal{M}$, we do not need to use the \emph{entire} category $\Comp^{\mathrm{op}}_{\mathcal{M}}$. We will return to this point in \S 8.
\end{remark}

We begin by observing that every ultracategory $\mathcal{M}$ can be identified with the full subcategory of $\Comp^{\mathrm{op}}_{\mathcal{M}}$ spanned by those pairs $(X, \mathcal{O}_X)$, where $X$ consists of a single point. To see this, we need the following general fact about left ultrafunctors:

\begin{proposition}[4.2.3]
Let $X = \{x\}$ be a one-point space. Then, for every ultracategory $\mathcal{M}$, the evaluation functor
\[
\FunLUlt(X, \mathcal{M}) \to \mathcal{M} \qquad F \mapsto F(x)
\]
is an equivalence of categories.
\end{proposition}

We postpone the proof of Proposition 4.2.3 for the moment.

\begin{remark}[4.2.4]
Let $\mathcal{M}$ be an ultracategory. Proposition 4.2.3 implies that for every object $M \in \mathcal{M}$, there is a unique left ultrastructure $\{\sigma_\mu\}$ on the functor
\[
F : X = \{x\} \to \mathcal{M} \qquad F(x) = M.
\]
Concretely, this left ultrastructure associates to each set $S$ and each ultrafilter $\mu$ on $S$ a map
\[
M = F\Big(\int_S x \, d\mu\Big) \xrightarrow{\ \sigma_\mu\ } \int_S F(x) \, d\mu = M^{\mu},
\]
which is given by the ultrapower diagonal of Example 1.3.4. Beware that this morphism is generally not invertible. In other words, there is generally no \emph{ultrafunctor} $\{x\} \to \mathcal{M}$ taking the value $M$. For example, if $\mathcal{M} = \mathbf{Set}$ is the category of sets, then the left ultrafunctor $F$ is an ultrafunctor if and only if $M$ is finite.
\end{remark}

\begin{notation}[4.2.5]
Let $\mathcal{M}$ be an ultracategory. For each object $M \in \mathcal{M}$, we let $\underline{M}$ denote the object of $\Comp_{\mathcal{M}}$ given by the pair $(\ast, \mathcal{O})$, where $\ast$ denotes the one-point space and $\mathcal{O} : \ast \to \mathcal{M}$ is the unique left ultrafunctor taking the value $M$.
\end{notation}

%% --- page 48 ---

\begin{remark}[4.2.6]
Let $\mathcal{M}$ be an ultracategory, let $M$ be an object of $\mathcal{M}$, and let $(X, \mathcal{O}_X)$ be any object of $\Comp_{\mathcal{M}}$. Every morphism $\underline{M} \to (X, \mathcal{O}_X)$ in the category $\Comp_{\mathcal{M}}$ determines a point $x \in X$ together with a morphism $\mathcal{O}_{X,x} \to M$ in the category $\mathcal{M}$. It follows from Proposition 4.2.3 that this construction induces a bijection
\[
\Hom_{\Comp_{\mathcal{M}}}\big(\underline{M}, (X, \mathcal{O}_X)\big) \to \coprod_{x \in X} \Hom_{\mathcal{M}}(\mathcal{O}_{X,x}, M).
\]
In particular, the construction $M \mapsto \underline{M}$ determines a fully faithful embedding $\mathcal{M} \to \Comp^{\mathrm{op}}_{\mathcal{M}}$.
\end{remark}

We can now state the main result of this section.

\begin{theorem}[4.2.7]
Let $\mathcal{M}$ be an ultracategory and let $\mathcal{M}' \subseteq \Comp^{\mathrm{op}}_{\mathcal{M}}$ denote the full subcategory spanned by those objects $(X, \mathcal{O}_X)$, where $X$ consists of a single point. Then:
\begin{enumerate}
\item[(a)] The category $\mathcal{M}'$ has ultraproducts in $\Comp^{\mathrm{op}}_{\mathcal{M}}$, in the sense of Construction 1.2.2).
\item[(b)] The construction $M \mapsto \underline{M}$ can be promoted to an equivalence of ultracategories $\mathcal{M} \to \mathcal{M}'$ (where $\mathcal{M}'$ is equipped with the ultrastructure of Proposition 1.3.7).
\end{enumerate}
\end{theorem}

The main ingredient in the proof of Theorem 4.2.7 is the fact that objects of the form $\underline{M}$ admit products in $\Comp^{\mathrm{op}}_{\mathcal{M}}$ (or, equivalently, coproducts in the category $\Comp_{\mathcal{M}}$). We now construct these.

\begin{proposition}[4.2.8]
Let $\mathcal{M}$ be an ultracategory and let $\{M_t\}_{t \in T}$ be a collection of objects of $\mathcal{M}$ indexed by a set $T$. Let $\mathcal{O}_{\beta T} : \beta T \to \mathcal{M}$ be the functor given by $\mathcal{O}_{\beta T, \nu} = \int_T M_t d\nu$. Then $\mathcal{O}_{\beta T}$ admits a left ultrastructure $\{\sigma_\mu\}$, which assigns to each map of sets $S \to \beta T$ given by a family of ultrafilters $\nu_\bullet = \{\nu_s\}_{s \in S}$ and each ultrafilter $\mu$ on $S$ the map
\[
\sigma_\mu : \mathcal{O}_{\beta T, \int_S \nu_s d\mu} = \int_T M_t d\Big(\int_S \nu_s d\mu\Big) \xrightarrow{\ \Delta_{\mu, \nu_\bullet}\ } \int_S \Big(\int_T M_t d\nu_s\Big) d\mu = \int_S \mathcal{O}_{\beta T, \nu_s} d\mu
\]
determined by the \emph{Fubini transformation} $\Delta_{\mu, \nu_\bullet}$.
\end{proposition}

\begin{proof}
We must show that the maps $\{\sigma_\mu\}$ satisfy conditions (0), (1), and (2) of Definition 1.4.1. Condition (0) is vacuous (since $X$ has only identity morphisms). To verify condition (1), suppose we are given a collection of ultrafilters $\nu_\bullet = \{\nu_s\}_{s \in S}$ on the set $T$ and an element $s_0 \in S$; we wish to show that the diagram
%% [NOTE: commutative diagram; every node, arrow and label is verbatim, geometric layout approximated.]
\[
\begin{array}{ccccc}
\int_T M_t d\Big(\int_S M_s d\delta_{s_0}\Big) & \xrightarrow{\ \ \Delta_{\delta_{s_0}, \nu_\bullet}\ } & \int_S \Big(\int_T N_t d\nu_s\Big) d\delta_{s_0} & & \\[6pt]
& \searrow{\ \ \text{(double-shaft arrow, no label)}} & & \swarrow{\ \ \epsilon_{S, s_0}} & \\[6pt]
& & \int_T M_t d\nu_{s_0} & & \\
\end{array}
\]
commutes in the category $\mathcal{M}$, which is axiom $(A)$ of Definition 1.3.1. To verify condition (2), we must show that for every collection $\nu_\bullet = \{\nu_s\}_{s \in S}$ of ultrafilters on $T$, every collection $\mu_\bullet = \{\mu_r\}_{r \in R}$ of ultrafilters on $S$, and every ultrafilter $\lambda$ on $R$, we have a commutative diagram
%% [NOTE: commutative diagram; every node, arrow and label is verbatim, geometric layout approximated.
%%  The two vertical bars "||" on the left are the source's own double line (an unlabelled identity/2-cell).]
\[
\begin{array}{ccccc}
\int_T M_t d\Big(\int_S \nu_s d\Big(\int_R \mu_r d\lambda\Big)\Big) & \xrightarrow{\ \ \Delta_{\int_R \mu_r d\lambda,\ \nu_\bullet}\ } & \int_S \Big(\int_T M_t d\nu_s\Big) d\Big(\int_R \mu_r d\lambda\Big) & & \\[6pt]
\Big\Vert & & & \Big\downarrow{\ \ \Delta_{\lambda, \mu_\bullet}} & \\[6pt]
\int_T M_t d\Big(\int_R \Big(\int_S \nu_s d\mu_r\Big) d\lambda\Big) & \xrightarrow{\ \ \Delta_{\lambda,\ \int_S \nu_s d\mu_\bullet}\ } & \int_R \Big(\int_T M_t d\Big(\int_S \nu_s d\mu_r\Big)\Big) d\lambda & \xrightarrow{\ \ \int_R \Delta_{\mu_r, \nu_\bullet} d\lambda\ } & \int_R \Big(\int_S \Big(\int_T N_t d\nu_s\Big) d\mu_r\Big) d\lambda, \\
\end{array}
\]
which is axiom $(C)$ of Definition 1.3.1.
%% [NOTE: printed glyph deviation, verified against a 400 dpi crop of page 48 in which both letters occur in
%%  one and the same formula: the middle node above is printed "int_T M_t d(int_S nu_s d mu_r)" and the
%%  right-hand node is printed "int_R(int_S(int_T N_t d nu_s)d mu_r)d lambda". That is, the source prints
%%  "N_t" in the position where the surrounding mathematics requires the object family "M_t" introduced in
%%  the statement of Proposition 4.2.8. The letter is a genuine italic "N" (with diagonal), visibly different
%%  from the italic "M" printed in the adjacent node; it is not a rendering artefact of this transcription.
%%  Reproduced as printed, NOT silently corrected. The same deviation occurs in the condition-(1) diagram
%%  above ("int_S(int_T N_t d nu_s)d delta_{s_0}") and in the proof of Proposition 4.2.9 on page 49.]
\end{proof}

In the situation of Proposition 4.2.8, the left ultrafunctor $\mathcal{O}_{\beta T}$ has a universal property.

\begin{proposition}[4.2.9]
Let $\mathcal{M}$ be an ultracategory, let $\{M_t\}_{t \in T}$ be a collection of objects of $\mathcal{M}$, and let $\mathcal{O}_{\beta T} : \beta T \to \mathcal{M}$ be the left ultrafunctor of Proposition 4.2.8. Let $\mathscr{F} : \beta T \to \mathcal{M}$ be any left ultrafunctor. For any natural transformation of left ultrafunctors $\alpha : \mathscr{F} \to \mathcal{O}_{\beta T}$ and any element $t \in T$, let $\alpha_t$ denote the composite map
\[
\mathscr{F}_{\delta_t} \xrightarrow{\ \alpha\ } \mathcal{O}_{\beta T}(\delta_t) = \int_T M_t d\delta_t \xrightarrow{\ \epsilon_{T,t}\ } M_t.
\]

%% --- page 49 ---

Then the construction $\alpha \mapsto \{\alpha_t\}_{t \in T}$ induces a bijection
\[
\theta : \Hom_{\FunLUlt(\beta T, \mathcal{M})}(\mathscr{F}, \mathcal{O}_{\beta T}) \to \prod_{t \in T} \Hom_{\mathcal{M}}(\mathscr{F}_{\delta_t}, M_t).
\]
\end{proposition}

\begin{proof}
Let $\{\sigma_\mu\}$ denote the left ultrastructure on $\mathcal{O}_{\beta T}$ constructed in Proposition 4.2.8 and let $\{\sigma'_\mu\}$ denote the left ultrastructure on $\mathscr{F}$. Let $\alpha : \mathscr{F} \to \mathcal{O}_{\beta T}$ be a natural transformation of left ultrafunctors. For each ultrafilter $\nu$ on $T$, we can write $\nu = \int_T \delta_t d\nu$, so that we have a commutative diagram
%% [NOTE: commutative diagram; every node, arrow and label is verbatim, geometric layout approximated.]
\[
\begin{array}{ccccc}
\mathscr{F}_{\int_T \delta_t d\nu} & \xrightarrow{\ \ \sigma'_\nu\ } & \int_T \mathscr{F}_{\delta_t} d\nu & & \\[6pt]
\Big\downarrow{\ \ \alpha(\nu)} & \swarrow{\ \ \int_T \alpha(\delta_t) d\nu} & \Big\downarrow{\ \ \int_T \alpha(\delta_t) d\nu} & \searrow{\ \ \int_T \alpha_t d\nu} & \\[6pt]
\mathcal{O}_{\beta T, \nu} & \xrightarrow{\ \ \sigma_\nu\ } & \int_T \mathcal{O}_{\beta T, \delta_t} d\nu & \xrightarrow{\ \ \int_T \epsilon_{T,t} d\nu\ } & \int_T N_t d\nu. \\
\end{array}
\]
It follows from Corollary 1.3.6 that the composition of the lower horizontal maps is the identity from $\int_T N_t d\nu$ to itself.
%% [NOTE: printed as "N_t" (see the note on page 48); context requires the family M_t of Proposition 4.2.9.
%%  Reproduced as printed.]
Consequently, $\alpha(\nu)$ is given by the composition
\[
\mathscr{F}_\nu = \mathscr{F}_{\int_T \delta_t d\nu} \xrightarrow{\ \sigma'_\nu\ } \int_T \mathscr{F}_{\delta_t} d\nu \xrightarrow{\ \int_T \alpha_t d\nu\ } \int_T N_t d\nu = \mathcal{O}_{\beta T, \nu}
\]
and is therefore completely determined by the maps $\{\alpha_t\}_{t \in T}$. This proves that $\theta$ is injective.

To show that $\theta$ is surjective, suppose we are given any collection of morphisms $\{f_t : \mathscr{F}_{\delta_t} \to M_t\}_{t \in T}$ in the category $\mathcal{M}$. Let $\alpha : \mathscr{F} \to \mathcal{O}_{\beta T}$ be the natural transformation which assigns to each ultrafilter $\nu$ on $T$ the composite map
\[
\mathscr{F}_\nu = \mathscr{F}_{\int_T \delta_t d\nu} \xrightarrow{\ \sigma'_\nu\ } \int_T \mathscr{F}_{\delta_t} d\nu \xrightarrow{\ \int_T f_t d\nu\ } \int_T N_t d\nu = \mathcal{O}_{\beta T, \nu}.
\]
%% [NOTE: the right-hand "N_t" is printed exactly as shown in the 400 dpi crop; context requires M_t.]
It follows from axiom (1) of Definition 1.4.1 that $\alpha_t = f_t$ for each $t \in T$. Consequently, to show that $\{f_t\}_{t \in T}$ belongs to the image of $\theta$, it will suffice to verify that $\alpha$ is a natural transformation of left ultrafunctors. Suppose we are given a collection of ultrafilters $\{\nu_s\}_{s \in S}$ on $T$ and an ultrafilter $\mu$ on $S$; we wish to show that the diagram
%% [NOTE: commutative diagram; every node, arrow and label is verbatim, geometric layout approximated.
%%  The "||" symbols are the source's own double lines (unlabelled identities).]
\[
\begin{array}{ccccc}
\mathscr{F}_{\int_S \nu_s d\mu} & \xrightarrow{\ \ \gamma'_\mu\ } & \int_S \mathscr{F}_{\nu_s} d\mu & & \\[6pt]
\Big\Vert & & \Big\Vert & & \\[6pt]
\mathscr{F}_{\int_T \delta_t d(\int_S \nu_s d\mu)} & & \int_S \mathscr{F}_{\int_T \delta_t d\nu_s} d\mu & & \\[6pt]
\Big\downarrow{\ \ \gamma'_{\int_S \nu_s d\mu}} & & \Big\downarrow{\ \ \int_S \gamma'_{\nu_s} d\mu} & & \\[6pt]
\int_T \mathscr{F}_{\delta_t} d\Big(\int_S \nu_s d\mu\Big) & \xrightarrow{\ \ \Delta_{\mu, \nu_\bullet}\ } & \int_S \Big(\int_T \mathscr{F}_{\delta_t} d\nu_s\Big) d\mu & & \\[6pt]
\Big\downarrow{\ \ \int_T f_t d(\int_S \nu_s d\mu)} & & \Big\downarrow{\ \ \int_S (\int_T f_t d\nu_s) d\mu} & & \\[6pt]
\int_T N_t d\Big(\int_S \nu_s d\mu\Big) & \xrightarrow{\ \ \Delta_{\mu, \nu_\bullet}\ } & \int_S \Big(\int_T N_t d\nu_s\Big) d\mu & & \\
\end{array}
\]
commutes. The commutativity of the upper rectangle follows from our assumption that $\{\gamma'_\mu\}$ is a left ultrastructure, and the commutativity of the lower square from the naturality of the Fubini transformations in the ultracategory $\mathcal{M}$.
%% [NOTE: the two occurrences of "N_t" in the bottom row are printed exactly as shown; context requires M_t.]
\end{proof}

%% --- page 50 ---

\begin{proof}[Proof of Proposition 4.2.3]
Let $X = \{x\}$ be a single point, which we identify with the Stone-\v{C}ech compactification of itself. We first note that for every object $M \in \mathcal{M}$, Proposition 4.2.8 produces a left ultrafunctor $\mathcal{O} : X \to \mathcal{M}$ satisfying $\mathcal{O}(x) = \int_X M\delta_x$, so that there exists an isomorphism $\epsilon_{X,x} : \mathcal{O}(x) \simeq M$. This proves the essential surjectivity of the evaluation map $\FunLUlt(X, \mathcal{M}) \to \mathcal{M}$. We will complete the proof by showing that it is fully faithful: that is, for every pair of left ultrafunctors $F, G : X \to \mathcal{M}$, the canonical map
\[
\Hom_{\FunLUlt(X, \mathcal{M})}(F, G) \to \Hom_{\mathcal{M}}(F(x), G(x))
\]
is bijective. Set $M = G(x)$ and define $\mathcal{O}$ as above. It follows from Proposition 4.2.9 that the isomorphism $\epsilon^{-1}_{X,x} : G(x) \simeq \mathcal{O}(x)$ lifts uniquely to a natural transformation of left ultrafunctors $\alpha : G \to \mathcal{O}$, which is automatically an isomorphism. We may therefore assume without loss of generality that $G = \mathcal{O}$, in which case the desired result follows from Proposition 4.2.9.
\end{proof}

%% [NOTE: the proof above is the complete proof of Proposition 4.2.3: four sentences, every substantive step
%%  of which is a citation of Proposition 4.2.8 or Proposition 4.2.9. It contains no independent argument.]

\begin{proposition}[4.2.10]
Let $\mathcal{M}$ be an ultracategory, let $\{M_t\}_{t \in T}$ be a collection of objects of $\mathcal{M}$, and let $\mathcal{O}_{\beta T} : \beta T \to \mathcal{M}$ be the left ultrafunctor of Proposition 4.2.8. For each $t \in T$, let $e_t : \underline{M_t} \to (\beta T, \mathcal{O}_{\beta T})$ denote the morphism in $\Comp_{\mathcal{M}}$ corresponding to the point $\delta_t \in \beta T$ and the isomorphism
\[
\mathcal{O}_{\beta T, \delta_t} = \int_T M_t d\delta_t \xrightarrow{\ \epsilon_{T,t}\ } M_t.
\]
Then the maps $e_t$ exhibit $(\beta T, \mathcal{O}_{\beta T})$ as a coproduct of the family of objects $\{\underline{M_t}\}_{t \in T}$ in the category $\Comp_{\mathcal{M}}$.
\end{proposition}

\begin{proof}
Combine Proposition 4.2.9 with Proposition 3.2.7.
\end{proof}

Proof of Theorem 4.2.7. Let $\{M_t\}_{t \in T}$ be a collection of objects of $\mathcal{M}$ and let $\mathcal{O}_{\beta T}$ be as in Proposition 4.2.8, so that Proposition 4.2.10 identifies $(\beta T, \mathcal{O}_{\beta T})$ with the product of $\{M_t\}_{t \in T_0}$ in $\Comp^{\mathrm{op}}_{\mathcal{M}}$.
%% [NOTE: the proof of Theorem 4.2.7 continues beyond this point onto printed page 51, which was not rendered;
%%  transcription stops here. It is included only to mark where the target item's page ends.]

\end{document}
